% !TEX program = pdflatex
\documentclass[10pt]{article}

\usepackage[utf8]{inputenc}
\usepackage[T1]{fontenc}
\usepackage[english]{babel}
\usepackage{lmodern}
\usepackage{microtype}
\usepackage[margin=0.72in]{geometry}
\usepackage{amsmath,amssymb,amsthm,mathtools}
\usepackage{booktabs}
\usepackage{xcolor}
\usepackage{enumitem}
\usepackage[most]{tcolorbox}
\usepackage{fancyhdr}
\usepackage[hidelinks]{hyperref}

\definecolor{CourseBlue}{HTML}{1A3C68}
\definecolor{CourseGreen}{HTML}{216E39}
\definecolor{CourseGold}{HTML}{B8860B}
\definecolor{CourseRed}{HTML}{9A2C2C}
\definecolor{SoftBlue}{HTML}{F2F6FA}
\definecolor{SoftGreen}{HTML}{F1F8F3}
\definecolor{SoftGold}{HTML}{FFF8E8}
\definecolor{SoftRed}{HTML}{FFF3F3}

\setlength{\parindent}{0pt}
\setlength{\parskip}{0.32em}
\setlist[itemize]{leftmargin=1.5em,itemsep=0.12em,topsep=0.2em}
\setlist[enumerate]{leftmargin=1.7em,itemsep=0.12em,topsep=0.2em}
\allowdisplaybreaks

\pagestyle{fancy}
\fancyhf{}
\lhead{Stochastic Calculus and Option Pricing}
\rhead{Lecture 6 Notes}
\cfoot{\thepage}
\renewcommand{\headrulewidth}{0.3pt}

\newcommand{\E}{\mathbb E}
\newcommand{\Pp}{\mathbb P}
\newcommand{\Q}{\mathbb Q}
\newcommand{\R}{\mathbb R}
\newcommand{\1}{\mathbf 1}
\newcommand{\cF}{\mathcal F}
\newcommand{\dd}{\,\mathrm d}
\newcommand{\Ito}{\mathrm{It\hat{o}}}
\newcommand{\phiN}{\varphi}
\DeclareMathOperator{\Var}{Var}
\DeclareMathOperator{\Cov}{Cov}

\theoremstyle{definition}
\newtheorem{definition}{Definition}[section]
\newtheorem{proposition}[definition]{Proposition}
\newtheorem{theorem}[definition]{Theorem}
\newtheorem{lemma}[definition]{Lemma}
\newtheorem{corollary}[definition]{Corollary}
\newtheorem{remark}[definition]{Remark}

\newtcolorbox{boardexample}[1]{
  enhanced,breakable,colback=SoftBlue,colframe=CourseBlue,
  boxrule=0.55pt,arc=1mm,left=1.5mm,right=1.5mm,top=1mm,bottom=1mm,
  before skip=0.55em,after skip=0.55em,
  title={Example --- #1},fonttitle=\bfseries}
\newtcolorbox{studentexercise}[1]{
  enhanced,breakable,colback=SoftGold,colframe=CourseGold,
  boxrule=0.55pt,arc=1mm,left=1.5mm,right=1.5mm,top=1mm,bottom=1mm,
  before skip=0.55em,after skip=0.35em,
  title={Exercise --- #1},fonttitle=\bfseries}
\newtcolorbox{shortsolution}[1][]{
  enhanced,breakable,colback=SoftGreen,colframe=CourseGreen,
  boxrule=0.55pt,arc=1mm,left=1.5mm,right=1.5mm,top=1mm,bottom=1mm,
  before skip=0.25em,after skip=0.6em,
  title={Short solution#1},fonttitle=\bfseries}
\newtcolorbox{keyidea}[1]{
  enhanced,breakable,colback=white,colframe=CourseGreen,
  boxrule=0.7pt,arc=1mm,left=1.5mm,right=1.5mm,top=1mm,bottom=1mm,
  before skip=0.5em,after skip=0.5em,
  title={#1},fonttitle=\bfseries}
\newtcolorbox{warningbox}[1]{
  enhanced,breakable,colback=SoftRed,colframe=CourseRed,
  boxrule=0.65pt,arc=1mm,left=1.5mm,right=1.5mm,top=1mm,bottom=1mm,
  before skip=0.5em,after skip=0.5em,
  title={#1},fonttitle=\bfseries}
\newtcolorbox{assumptionbox}[1]{
  enhanced,breakable,colback=SoftGold,colframe=CourseGold,
  boxrule=0.65pt,arc=1mm,left=1.5mm,right=1.5mm,top=1mm,bottom=1mm,
  before skip=0.5em,after skip=0.5em,
  title={#1},fonttitle=\bfseries}

\title{\textbf{Girsanov and Risk-Neutral Pricing}\\
\large Lecture 6 Notes --- Densities, Drift Changes, and Risk-Neutral Dynamics}
\author{}
\date{}

\begin{document}
\maketitle
\vspace{-1.4em}

\begin{keyidea}{Route through Lecture 6}
A rare Gaussian tail first motivates reweighting. Finite scenarios and
Gaussian tilting explain the density; exponential weights extend the idea
to Brownian paths. Girsanov identifies the new
Brownian motion, and the Black--Scholes example chooses its shift so that
the discounted stock has no drift. Lecture 7 starts from this result and
develops replication, conditional-expectation pricing, and completeness.
\end{keyidea}

% ============================================================
\section{Densities and Bayes' Rule}
% ============================================================

\subsection{A rare Gaussian tail motivates reweighting}

Let \(X\sim N(0,1)\) under \(\Pp\), with \(\Phi\) the standard normal CDF, and consider
\[
p=\Pp(X>5)=1-\Phi(5)\approx2.87\times10^{-7}.
\]
The direct Monte Carlo estimate from \(n\) independent draws is
\(\widehat p_{\mathrm{MC}}=n^{-1}\sum_{i=1}^n\1_{\{X_i>5\}}\).
Even for \(n=10^6\), the expected number of tail observations is only
\(np\approx0.29\). The probability of seeing none is
\((1-p)^n\approx0.75\), in which case the estimate is zero.

To visit the tail more often, sample under a new law \(\Q\) for which the
same variable has distribution \(X\sim N(5,1)\). Then
\(\Q(X>5)=1/2\). This \(\Q\) is chosen for simulation; it need not be a
pricing measure. The likelihood correction that recovers \(p\) is derived
in the Gaussian tilting section below.

\subsection{Finite spaces reveal the mechanism}

Let \(\Omega=\{\omega_1,\ldots,\omega_m\}\), and suppose
\(\Pp(\omega_i)>0\) and \(\Q(\omega_i)>0\). Define the density ratio
\[
Z(\omega_i)=\frac{\Q(\omega_i)}{\Pp(\omega_i)}.
\]
Then
\[
Z>0,
\qquad
\E^{\Pp}[Z]=1,
\qquad
\boxed{\E^{\Q}[Y]=\E^{\Pp}[ZY].}
\]
A state with \(Z>1\) is upweighted and one with \(Z<1\) is downweighted.

\begin{boardexample}{a three-state reweighting}
Consider
\[
\begin{array}{c|ccc}
 & \omega_1 & \omega_2 & \omega_3\\ \hline
\Pp & 1/2 & 1/3 & 1/6\\
\Q  & 1/4 & 1/2 & 1/4
\end{array}
\]
Then
\[
Z(\omega_1)=\frac12,
\qquad
Z(\omega_2)=Z(\omega_3)=\frac32,
\]
and indeed
\[
\E^{\Pp}[Z]
=\frac12\frac12+\frac13\frac32+\frac16\frac32=1.
\]
For \(Y=(-2,1,4)\), direct computation and reweighting agree:
\[
\E^{\Q}[Y]=1=\E^{\Pp}[ZY].
\]
\end{boardexample}

\subsection{Radon--Nikodym densities and equivalence}

\begin{definition}[Change of measure by a density]
Let \(Z\geq0\) be \(\cF_T\)-measurable with \(\E^{\Pp}[Z]=1\). Define
\[
\Q(A)=\E^{\Pp}[Z\1_A],
\qquad A\in\cF_T.
\]
Then \(\Q\) is a probability measure and
\[
Z=\frac{\dd\Q}{\dd\Pp}\bigg|_{\cF_T}.
\]
For every \(\Q\)-integrable \(Y\),
\[
\boxed{\E^{\Q}[Y]=\E^{\Pp}[ZY].}
\]
\end{definition}

Absolute continuity \(\Q\ll\Pp\) means that a \(\Pp\)-null event remains
impossible under \(\Q\). If \(Z>0\) \(\Pp\)-almost surely, then
\(\Pp\sim\Q\): both measures have exactly the same null events. In that case
\[
\frac{\dd\Pp}{\dd\Q}=\frac1Z
\quad\Q\text{-a.s.},
\qquad
\E^{\Q}\!\left[\frac1Z\right]=1.
\]

\begin{warningbox}{Three properties with different roles}
\begin{itemize}
  \item \(Z\geq0\) makes \(\Q(A)=\E^{\Pp}[Z\1_A]\) nonnegative.
  \item \(\E^{\Pp}[Z]=1\) gives total mass one.
  \item \(Z>0\) gives equivalence, which prevents the change of measure from
        deleting possible scenarios.
\end{itemize}
Do not replace one property by another.
\end{warningbox}

\subsection{Conditional normalization and the density process}

Let \(\mathcal G\subseteq\cF_T\). If the denominator is positive, Bayes'
rule under a change of measure is
\[
\boxed{
\E^{\Q}[Y\mid\mathcal G]
=\frac{\E^{\Pp}[ZY\mid\mathcal G]}
       {\E^{\Pp}[Z\mid\mathcal G]}.}
\]
The numerator reweights future outcomes. The denominator renormalizes those
weights after the information in \(\mathcal G\) is known.

To prove the formula, set
\[
X=\frac{\E^{\Pp}[ZY\mid\mathcal G]}
        {\E^{\Pp}[Z\mid\mathcal G]}.
\]
The variable \(X\) is \(\mathcal G\)-measurable. For every
\(A\in\mathcal G\),
\[
\begin{aligned}
\E^{\Q}[\1_AX]
&=\E^{\Pp}[Z\1_AX]\\
&=\E^{\Pp}[\1_AX\E^{\Pp}[Z\mid\mathcal G]]\\
&=\E^{\Pp}[\1_A\E^{\Pp}[ZY\mid\mathcal G]]\\
&=\E^{\Pp}[\1_AZY]
=\E^{\Q}[\1_AY].
\end{aligned}
\]
This is the defining test for \(X=\E^{\Q}[Y\mid\mathcal G]\).

Given a terminal density \(Z_T\), define
\[
\boxed{Z_t=\E^{\Pp}[Z_T\mid\cF_t].}
\]
Then \(Z=(Z_t)_{0\leq t\leq T}\) is the density process, and
\[
\boxed{
\E^{\Q}[Y\mid\cF_t]
=\frac1{Z_t}\E^{\Pp}[Z_TY\mid\cF_t].}
\]

\begin{keyidea}{Information is not changed by reweighting}
If \(Y\) is already \(\mathcal G\)-measurable, then
\[
\E^{\Pp}[ZY\mid\mathcal G]
=Y\E^{\Pp}[Z\mid\mathcal G],
\]
so \(\E^{\Q}[Y\mid\mathcal G]=Y\). Changing measure changes probabilities,
not what has already been observed.
\end{keyidea}

% ============================================================
\section{Gaussian Tilting: the Finite-Dimensional Prototype}
% ============================================================

\subsection{Completing the square fixes the sign}

Let \(X\sim N(0,1)\) under \(\Pp\), and define
\[
Z=\exp\!\left(-\theta X-\frac12\theta^2\right).
\]
The Gaussian moment-generating function gives
\[
\E^{\Pp}[Z]
=e^{-\theta^2/2}\E^{\Pp}[e^{-\theta X}]=1.
\]
Under \(\Q\), the density of \(X\) is
\[
e^{-\theta x-\theta^2/2}\frac{e^{-x^2/2}}{\sqrt{2\pi}}
=\frac1{\sqrt{2\pi}}e^{-(x+\theta)^2/2}.
\]
Therefore
\[
\boxed{X\sim N(-\theta,1)\text{ under }\Q.}
\]
For \(\theta>0\), the factor \(e^{-\theta X}\) favors negative values. This
quick qualitative check is often enough to detect a sign error.

More generally, if \(X\sim N(0,\sigma^2)\) and we want mean \(m\) under
\(\Q\), take
\[
\boxed{
Z=\exp\!\left(\frac{m}{\sigma^2}X
-\frac{m^2}{2\sigma^2}\right).}
\]
Then \(X\sim N(m,\sigma^2)\) under \(\Q\). The mean changes; the variance
does not.

\begin{studentexercise}{reverse-engineer the normalizer}
Suppose \(X\sim N(0,v)\) and \(Z=e^{aX+b}\). Find the unique \(b\) for which
\(\E[Z]=1\), then identify the law of \(X\) under \(\dd\Q=Z\dd\Pp\).
\end{studentexercise}
\begin{shortsolution}
Since \(\E[e^{aX}]=e^{a^2v/2}\), normalization requires
\(b=-a^2v/2\). Completing the square gives
\[
\boxed{X\sim N(av,v)\text{ under }\Q.}
\]
The linear coefficient controls the new mean; the quadratic compensator
ensures total mass one.
\end{shortsolution}

\subsection{Reweighting and translation}

For every bounded measurable \(f\),
\[
\boxed{
\E^{\Q}[f(X)]
=\E^{\Pp}[Zf(X)]
=\E^{\Pp}[f(X-\theta)]}
\]
in the standard-normal example above. The first equality is the definition
of the new measure. The second is a Gaussian distributional identity. It is
important not to confuse these two logical steps.

\begin{keyidea}{What survives in continuous time}
A Brownian path is an infinite collection of compatible Gaussian variables.
The continuous-time density will again contain
\[
\text{linear Gaussian exposure}
-\frac12\text{variance of that exposure}.
\]
This is the origin of the stochastic exponential.
\end{keyidea}

\subsection{A shifted tail with and without likelihood weights}

Return to \(p=\Pp(X>5)\) with \(X\sim N(0,1)\) under \(\Pp\). The Gaussian
tilt with target mean \(5\) gives \(X\sim N(5,1)\) under \(\Q\), with
\[
\frac{\dd\Q}{\dd\Pp}(x)=e^{5x-25/2},
\qquad
w(x)=\frac{\dd\Pp}{\dd\Q}(x)=e^{-5x+25/2}.
\]
For independent draws \(X_1,\ldots,X_n\) under \(\Q\), merely shifting
the sampling mean gives
\[
\frac1n\sum_{i=1}^n\1_{\{X_i>5\}}
\longrightarrow\Q(X>5)=\frac12.
\]
This is the tail probability under the new law. To estimate the original
probability, reweight each observation:
\[
\boxed{
\widehat p_{\mathrm{IS}}
=\frac1n\sum_{i=1}^n w(X_i)\1_{\{X_i>5\}},
\qquad
\E^{\Q}[\widehat p_{\mathrm{IS}}]
=\E^{\Q}[w(X)\1_{\{X>5\}}]
=\Pp(X>5)=p.}
\]
This is importance sampling: the shift makes tail observations frequent,
and the inverse density restores their original probability weight.

% ============================================================
\section{Why Change the Probability Measure?}
% ============================================================

\subsection{One random payoff, two systems of weights}

Let \(H:\Omega\to\R\) be a terminal payoff. A probability measure assigns
weights to the scenarios \(\omega\in\Omega\); it does not alter the value
\(H(\omega)\). If \(\Pp\) is replaced by an equivalent measure \(\Q\), then
the sample space, filtration, stock paths, and payoff remain the same.
Only their probabilities change.

\begin{keyidea}{Physical and pricing viewpoints}
\begin{itemize}
  \item Under \(\Pp\), parameters are interpreted statistically: they describe
        frequencies or forecasts in the model.
  \item Under \(\Q\), probabilities are selected so that prices expressed in
        units of a chosen numeraire are martingales.
  \item A pricing measure is therefore constrained by traded prices. It is not
        a claim about investor beliefs or the most likely real-world outcome.
\end{itemize}
\end{keyidea}

Suppose, for the moment, that the stock and bank account satisfy under
\(\Pp\)
\[
\frac{\dd S_t}{S_t}=\mu\dd t+\sigma\dd W_t,
\qquad
\dd B_t=rB_t\dd t,
\qquad B_t=e^{rt}.
\]
The discounted stock is \(\bar S_t=S_t/B_t=e^{-rt}S_t\). Since the discount
factor has finite variation, the product rule gives
\[
\boxed{
\frac{\dd\bar S_t}{\bar S_t}
=(\mu-r)\dd t+\sigma\dd W_t.}
\]
Thus the term to remove is the excess drift \(\mu-r\). Measured per unit of
volatility, it is
\[
\lambda=\frac{\mu-r}{\sigma},
\]
the market price of Brownian risk in this one-factor model.

\begin{studentexercise}{discount before changing measure}
Let \(B_t=\exp(\int_0^t r_s\dd s)\) and
\(\dd S_t/S_t=\mu_t\dd t+\sigma_t\dd W_t\). Derive the SDE of
\(S_t/B_t\) without assuming constant coefficients.
\end{studentexercise}
\begin{shortsolution}
Writing \(D_t=B_t^{-1}\), we have \(\dd D_t=-r_tD_t\dd t\) and
\([D,S]=0\). Hence
\[
\boxed{
\frac{\dd(D_tS_t)}{D_tS_t}
=(\mu_t-r_t)\dd t+\sigma_t\dd W_t.}
\]
Discounting identifies exactly which drift Girsanov must eliminate.
\end{shortsolution}

\begin{warningbox}{Expectation under \(\Pp\) is not an arbitrage-free price}
In general,
\[
e^{-r(T-t)}\E^{\Pp}[H\mid\cF_t]
\]
is neither a replication price nor a no-arbitrage price. The physical drift
and risk preferences can enter that expectation. Pricing by expectation only
becomes valid after specifying the appropriate pricing measure and numeraire.
\end{warningbox}

% ============================================================
\section{From Gaussian Tilting to Brownian Paths}
% ============================================================

\subsection{A deterministic path exposure}

Let \(W\) be Brownian motion under \(\Pp\), and let
\(h\in L^2([0,T])\) be deterministic. The stochastic integral
\[
X=\int_0^Th_s\dd W_s
\]
is centered Gaussian with variance
\[
v=\int_0^Th_s^2\dd s.
\]
Consequently,
\[
\boxed{
Z_T=\exp\!\left(
\int_0^Th_s\dd W_s-\frac12\int_0^Th_s^2\dd s
\right)}
\]
satisfies \(\E^{\Pp}[Z_T]=1\).

\begin{theorem}[Cameron--Martin drift shift]
Define \(\Q\) on \(\cF_T\) by \(\dd\Q=Z_T\dd\Pp\). Then
\[
\boxed{B_t=W_t-\int_0^th_s\dd s}
\]
is Brownian motion under \(\Q\). Equivalently,
\[
\boxed{\dd W_t=\dd B_t+h_t\dd t.}
\]
\end{theorem}

The density rewards paths whose increments align with \(h\). Under the new
weights, \(W\) has accumulated drift \(\int_0^th_s\dd s\), and subtracting
that drift recovers a Brownian motion.

\begin{boardexample}{constant exposure}
If \(h_t\equiv\theta\), then
\[
Z_T=e^{\theta W_T-\theta^2T/2},
\qquad
B_t=W_t-\theta t.
\]
Thus \(W_T=B_T+\theta T\), and
\[
\boxed{W_T\sim N(\theta T,T)\text{ under }\Q.}
\]
This is precisely Gaussian tilting with a variance \(T\).
\end{boardexample}

\subsection{Why an adapted exposure is harder}

If \(h_t\) is replaced by a random progressively measurable process
\(\phi_t\), then \(\int_0^T\phi_s\dd W_s\) need not be Gaussian. Conditional
on the entire path of \(\phi\), the integral is not generally an independent
normal random variable, because \(\phi\) may depend on the same Brownian
motion. The single Gaussian-mgf argument therefore breaks down.

This is not a problem of notation. It is the reason we need stochastic
calculus and a genuine martingale condition.

% ============================================================
\section{The Stochastic Exponential and Its Validity}
% ============================================================

\subsection{The compensation term is forced by \texorpdfstring{\(\Ito\)}{Ito}}

Let \(\phi\) be progressively measurable and suppose
\[
\int_0^T\phi_s^2\dd s<\infty
\quad\text{almost surely}.
\]
Define
\[
M_t=\int_0^t\phi_s\dd W_s,
\qquad
[M]_t=\int_0^t\phi_s^2\dd s,
\]
and its stochastic exponential
\[
\boxed{
Z_t=\mathcal E(M)_t
=\exp\!\left(M_t-\frac12[M]_t\right).}
\]
Set \(Y_t=M_t-\tfrac12[M]_t\), so \(Z_t=e^{Y_t}\). Since
\[
\dd Y_t=\phi_t\dd W_t-\frac12\phi_t^2\dd t,
\qquad
\dd[Y]_t=\phi_t^2\dd t,
\]
\(\Ito\)'s formula gives
\[
\begin{aligned}
\dd Z_t
&=Z_t\dd Y_t+\frac12Z_t\dd[Y]_t\\
&=Z_t\phi_t\dd W_t
-\frac12Z_t\phi_t^2\dd t
+\frac12Z_t\phi_t^2\dd t.
\end{aligned}
\]
Therefore
\[
\boxed{\dd Z_t=Z_t\phi_t\dd W_t,\qquad Z_0=1.}
\]
The coefficient (1/2) is exactly what cancels the quadratic-variation
drift.

\begin{studentexercise}{detect the only valid compensator}
For \(c\in\R\), define
\[
Z_t^{(c)}=\exp\!\left(
\int_0^t\phi_s\dd W_s-c\int_0^t\phi_s^2\dd s
\right).
\]
Find the drift of \(Z^{(c)}\) and determine the unique \(c\) that makes it a
local martingale for every admissible \(\phi\).
\end{studentexercise}
\begin{shortsolution}
\(\Ito\)'s formula gives
\[
\frac{\dd Z_t^{(c)}}{Z_t^{(c)}}
=\phi_t\dd W_t+\left(\frac12-c\right)\phi_t^2\dd t.
\]
Hence
\[
\boxed{c=\frac12.}
\]
The correction is not a convention; it is forced by quadratic variation.
\end{shortsolution}

\subsection{Local martingale is not enough}

The stochastic exponential is positive and is a local martingale. Every
positive local martingale is a supermartingale, so
\[
\E^{\Pp}[Z_t]\leq Z_0=1.
\]
It can lose mass. If \(\E^{\Pp}[Z_T]<1\), the terminal variable cannot define
a probability measure of total mass one.

\begin{warningbox}{The most common unjustified step}
Writing \(\dd Z_t=Z_t\phi_t\dd W_t\) proves only that \(Z\) is a local
martingale. It does not prove \(\E[Z_T]=1\). Before defining \(\Q\), one must
verify a true-martingale condition or explicitly assume it.
\end{warningbox}

\begin{theorem}[Novikov's sufficient condition]
If
\[
\boxed{
\E^{\Pp}\!\left[
\exp\!\left(\frac12\int_0^T\phi_s^2\dd s\right)
\right]<\infty,}
\]
then \(Z=\mathcal E(\int\phi\dd W)\) is a true martingale on ([0,T]), and
\(\E^{\Pp}[Z_T]=1\).
\end{theorem}

If \(\lvert\phi_t\rvert\leq C\), then
\[
\int_0^T\phi_s^2\dd s\leq C^2T,
\]
and Novikov follows immediately. By contrast, almost-sure finiteness of
\(\int\phi^2\) does not automatically control its exponential moment.

\begin{boardexample}{a bounded adapted exposure}
Fix \(t_0<T\) and let \(\xi\) be bounded and \(\cF_{t_0}\)-measurable. Set
\[
\phi_t=\xi\1_{(t_0,T]}(t).
\]
This process is progressively measurable, and
\[
\int_0^T\phi_t^2\dd t=\xi^2(T-t_0)\leq
\|\xi\|_\infty^2T.
\]
Novikov holds even though the drift shift is random. Economically, the
future change of drift may depend on information already observed at time
\(t_0\), but not on future increments.
\end{boardexample}

\begin{keyidea}{Localization separates algebra from validity}
If \(\phi\) is unbounded, define stopping times
\[
\tau_n=\inf\left\{t:\int_0^t\phi_s^2\dd s\geq n\right\}\wedge T.
\]
The stopped integrand \(\phi_t\1_{\{t\leq\tau_n\}}\) has bounded accumulated
energy, so its stochastic exponential is a true martingale. This validates
the drift calculation up to \(\tau_n\). Passing to the unstopped limit still
requires a separate uniform-integrability or convergence argument.
\end{keyidea}

% ============================================================
\section{Girsanov Moves Drift into the Measure}
% ============================================================

\subsection{The theorem and its sign convention}

\begin{theorem}[Girsanov, one-dimensional version]
Let
\[
Z_t=\exp\!\left(
\int_0^t\phi_s\dd W_s
-\frac12\int_0^t\phi_s^2\dd s
\right)
\]
be a true \(\Pp\)-martingale, and define
\[
\frac{\dd\Q}{\dd\Pp}\bigg|_{\cF_T}=Z_T.
\]
Then
\[
\boxed{W_t^{\Q}=W_t-\int_0^t\phi_s\dd s}
\]
is Brownian motion under \(\Q\).
\end{theorem}

Always solve this identity for the old differential before substituting:
\[
\boxed{\dd W_t=\dd W_t^{\Q}+\phi_t\dd t.}
\]
Hence, if under \(\Pp\)
\[
\dd X_t=b_t\dd t+a_t\dd W_t,
\]
then under \(\Q\)
\[
\boxed{
\dd X_t=(b_t+a_t\phi_t)\dd t+a_t\dd W_t^{\Q}.}
\]
Only the drift changes. The diffusion coefficient, and therefore the local
quadratic variation, is unchanged:
\[
\dd[X]_t=a_t^2\dd t
\]
under both measures.

\begin{warningbox}{A sign table worth memorizing}
\[
\begin{array}{c|c|c}
\text{density exponent} & \text{new Brownian motion} & \text{new drift of }X\\ \hline
+\int\phi\dd W-\tfrac12\int\phi^2\dd t
& W^{\Q}=W-\int\phi\dd t & b+a\phi\\
-\int\lambda\dd W-\tfrac12\int\lambda^2\dd t
& W^{\Q}=W+\int\lambda\dd t & b-a\lambda
\end{array}
\]
The second row is the form most often used in asset pricing.
\end{warningbox}

\begin{boardexample}{remove a deterministic drift}
Suppose
\[
\dd X_t=\alpha_t\dd t+\beta_t\dd W_t,
\qquad \beta_t\neq0,
\]
with deterministic coefficients. To make \(X\) driftless, choose
\[
\phi_t=-\frac{\alpha_t}{\beta_t}.
\]
Then, provided the exponential is a true martingale,
\[
\dd X_t=\beta_t\dd W_t^{\Q}.
\]
This is first an algebraic candidate. Its validity still depends on the
density.
\end{boardexample}

\subsection{Write the density in each drift-change example}

For \(\dd X_t=\mu\dd t+\sigma\dd W_t\), \(\sigma>0\), target drift
\(\nu\) under \(\Q\) by taking
\[
\lambda=\frac{\mu-\nu}{\sigma},\qquad
Z_t=e^{-\lambda W_t-\lambda^2t/2},\qquad
\frac{\dd\Q}{\dd\Pp}\bigg|_{\cF_T}=Z_T.
\]
Then \(W_t^{\Q}=W_t+\lambda t\) and
\(\dd X_t=\nu\dd t+\sigma\dd W_t^{\Q}\).

For the adapted example
\(\dd X_t=\sin(W_t)\dd t+\dd W_t\), \(X_0=0\), set
\[
Z_t=\exp\!\left(-\int_0^t\sin(W_s)\dd W_s
-\frac12\int_0^t\sin^2(W_s)\dd s\right),\qquad
\frac{\dd\Q}{\dd\Pp}\bigg|_{\cF_T}=Z_T.
\]
The integrand is bounded, so Novikov applies. The process
\(W_t^{\Q}=W_t+\int_0^t\sin(W_s)\dd s\) is Brownian under \(\Q\), and
\(\dd X_t=\dd W_t^{\Q}\).

For a geometric diffusion \(\dd S_t/S_t=\mu\dd t+\sigma\dd W_t\), use the
same \(\lambda=(\mu-\nu)/\sigma\) and the full density
\[
Z_t=e^{-\lambda W_t-\lambda^2t/2},\qquad
\frac{\dd\Q}{\dd\Pp}\bigg|_{\cF_T}=Z_T.
\]
It yields \(\dd S_t/S_t=\nu\dd t+\sigma\dd W_t^{\Q}\). For the
money-market pricing measure, the target is \(\nu=r\).

\subsection{Conditional expectations and path-dependent claims}

For any \(Y\in L^1(\Q)\),
\[
\E^{\Q}[Y\mid\cF_t]
=\frac1{Z_t}\E^{\Pp}[Z_TY\mid\cF_t].
\]
This identity remains valid when \(Y\) depends on the whole path. For example,
\[
Y=\1_{\{\max_{0\leq u\leq T}W_u>K\}}
\]
does not become a function of \(W_T\) merely because the measure changes.
Girsanov changes the weight of each path, not the information contained in
the payoff.

In the geometric example with constant \(\nu,\sigma\), the two routes give
the same conditional mean:
\[
\boxed{\E^{\Q}[S_T\mid\cF_t]
=S_te^{\nu(T-t)}
=\frac1{Z_t}\E^{\Pp}[Z_TS_T\mid\cF_t].}
\]

\begin{keyidea}{Three checkpoints for every Girsanov problem}
\begin{enumerate}
  \item \textbf{Density:} write the linear stochastic integral and the
        negative half-energy compensator with the intended sign.
  \item \textbf{Validity:} prove that the exponential is a true martingale,
        or clearly state the assumption that guarantees it.
  \item \textbf{Substitution:} identify the new Brownian motion, solve for
        the old \(\dd W\), and only then rewrite the SDE.
\end{enumerate}
\end{keyidea}

% ============================================================
\section{The Risk-Neutral Measure in Black--Scholes}
% ============================================================

\subsection{Constructing the pricing density}

Under \(\Pp\), assume
\[
\frac{\dd S_t}{S_t}=\mu_t\dd t+\sigma_t\dd W_t,
\qquad
\dd B_t=r_tB_t\dd t,
\qquad \sigma_t\neq0.
\]
Define
\[
\lambda_t=\frac{\mu_t-r_t}{\sigma_t}.
\]
The pricing-density candidate is
\[
\boxed{
Z_t=\exp\!\left(
-\int_0^t\lambda_s\dd W_s
-\frac12\int_0^t\lambda_s^2\dd s
\right).}
\]
Assume it is a true martingale and set
\(\dd\Q/\dd\Pp|_{\cF_T}=Z_T\). Girsanov gives
\[
\boxed{
W_t^{\Q}=W_t+\int_0^t\lambda_s\dd s,
\qquad
\dd W_t=\dd W_t^{\Q}-\lambda_t\dd t.}
\]
Substituting into the stock SDE,
\[
\begin{aligned}
\frac{\dd S_t}{S_t}
&=\mu_t\dd t+\sigma_t
  (\dd W_t^{\Q}-\lambda_t\dd t)\\
&=(\mu_t-\sigma_t\lambda_t)\dd t
  +\sigma_t\dd W_t^{\Q}\\
&=\boxed{r_t\dd t+\sigma_t\dd W_t^{\Q}.}
\end{aligned}
\]
The statistical drift \(\mu_t\) has been replaced by the short rate, while
the volatility remains unchanged.

Let
\[
D_t=B_t^{-1}=\exp\!\left(-\int_0^tr_s\dd s\right).
\]
Under \(\Q\),
\[
\boxed{
\dd(D_tS_t)=\sigma_tD_tS_t\dd W_t^{\Q}.}
\]
Thus (DS) is a local martingale and, under suitable integrability, a true
martingale.

\subsection{A call formula from two Gaussian tails}

For constant \(r,\sigma\) and \(\tau=T-t\), the risk-neutral SDE gives
\[
S_T=S_t\exp\!\left((r-\tfrac12\sigma^2)\tau
+\sigma(W_T^{\Q}-W_t^{\Q})\right).
\]
Write
\[
d_-=\frac{\log(S_t/K)+(r-\frac12\sigma^2)\tau}
          {\sigma\sqrt\tau},\qquad
d_+=d_-+\sigma\sqrt\tau.
\]
The cash-or-nothing tail and the stock-weighted tail are
\[
\Q(S_T>K\mid\cF_t)=\Phi(d_-),\qquad
e^{-r\tau}\E^{\Q}[S_T\1_{\{S_T>K\}}\mid\cF_t]
=S_t\Phi(d_+).
\]
Since \((S_T-K)^+=S_T\1_{\{S_T>K\}}-K\1_{\{S_T>K\}}\),
\[
\boxed{C_t=S_t\Phi(d_+)-Ke^{-r\tau}\Phi(d_-),\qquad
\Delta_t=\Phi(d_+).}
\]
The pricing identity and the replication argument behind it are developed
in Lecture 7.

\begin{studentexercise}{catch the wrong pricing sign}
For constant \(\lambda=(\mu-r)/\sigma\), suppose one mistakenly uses
\[
\widetilde Z_T=e^{+\lambda W_T-\lambda^2T/2}.
\]
Find the stock drift under the corresponding measure.
\end{studentexercise}
\begin{shortsolution}
Now \(\widetilde W_t=W_t-\lambda t\) is Brownian, so
\(\dd W_t=\dd\widetilde W_t+\lambda\dd t\). Therefore
\[
\frac{\dd S_t}{S_t}
=(\mu+\sigma\lambda)\dd t+\sigma\dd\widetilde W_t
=(2\mu-r)\dd t+\sigma\dd\widetilde W_t.
\]
The discounted drift is \(2(\mu-r)\), so the sign is wrong.
\end{shortsolution}

\subsection{What ``risk-neutral'' does and does not mean}

\begin{definition}[Equivalent martingale measure]
Relative to the positive numeraire \(B\), a probability measure \(\Q\) is an
equivalent martingale measure if
\begin{enumerate}
  \item \(\Q\sim\Pp\) on \(\cF_T\), and
  \item every discounted traded price \(S^i/B\) is a \(\Q\)-martingale
        under the required admissibility and integrability conditions.
\end{enumerate}
\end{definition}

The term ``risk-neutral'' means that traded assets have drift \(r_t\) when
expressed in currency units, equivalently zero drift after discounting. It
does not mean:
\begin{itemize}
  \item that real investors are indifferent to risk;
  \item that \(\Q\) is the empirical distribution of future returns;
  \item that all payoffs have expected return \(r\) under \(\Pp\);
  \item that changing measure removes randomness.
\end{itemize}

\begin{keyidea}{Why the physical drift disappears from option prices}
Replication and Girsanov express the same no-arbitrage restriction in two
languages. Delta hedging cancels the stock's Brownian exposure and therefore
its associated physical drift. Girsanov reweights paths so that discounted
traded assets are martingales. Both routes produce the same pricing PDE and
the same value.
\end{keyidea}

% ============================================================
\section{Guided Practice and Common Traps}
% ============================================================

\subsection{A complete change-of-measure drill}

\begin{studentexercise}{a constant Brownian shift}
Let \(W\) be Brownian motion under \(\Pp\), fix \(a\in\R\), and define
\[
Z_T=e^{-aW_T-a^2T/2},
\qquad
\frac{\dd\Q}{\dd\Pp}=Z_T.
\]
\begin{enumerate}
  \item Verify \(\E^{\Pp}[Z_T]=1\).
  \item Identify the \(\Q\)-Brownian motion.
  \item Find the law of \(W_T\) under \(\Q\).
  \item Rewrite \(\dd X_t=b\dd t+c\dd W_t\) under \(\Q\).
\end{enumerate}
\end{studentexercise}
\begin{shortsolution}
The Gaussian mgf gives the normalization. Here \(\phi=-a\), so
\[
W_t^{\Q}=W_t+at,
\qquad
W_T=W_T^{\Q}-aT.
\]
Hence
\[
W_T\sim N(-aT,T)\quad\text{under }\Q,
\]
and \(\dd W_t=\dd W_t^{\Q}-a\dd t\) yields
\[
\boxed{\dd X_t=(b-ac)\dd t+c\dd W_t^{\Q}.}
\]
\end{shortsolution}

\subsection{Diagnose candidate densities}

\begin{studentexercise}{which conclusions are justified?}
Consider
\[
\begin{array}{c|c}
\text{candidate} & \text{available information}\\ \hline
e^{\theta W_T-\theta^2T/2} & \theta\in\R\\
e^{\int_0^T\phi_s\dd W_s-\frac12\int_0^T\phi_s^2\dd s}
& |\phi_s|\leq C\\
e^{\int_0^T\psi_s\dd W_s-\frac12\int_0^T\psi_s^2\dd s}
& \int_0^T\psi_s^2\dd s<\infty\text{ a.s.}
\end{array}
\]
For each row, decide whether the stated information is enough to define a
probability measure.
\end{studentexercise}
\begin{shortsolution}
The first density has expectation one by the Gaussian mgf. The second
satisfies Novikov because the integrand is bounded. The third is only known
to be a positive local martingale; finite energy almost surely does not imply
the exponential moment required by Novikov or any equivalent criterion.
\end{shortsolution}

\subsection{Error diagnosis}

\begin{studentexercise}{correct six statements}
Correct the following claims:
\begin{enumerate}
  \item A change of measure changes the realized payoff \(H(\omega)\).
  \item \(Z_T\geq0\) is enough to define a probability measure.
  \item The conditional change-of-measure formula has denominator \(Z_T\).
  \item Zero drift of a stochastic exponential proves it is a true martingale.
  \item Girsanov changes both drift and volatility.
  \item Under a risk-neutral measure, every undiscounted traded price is a
        martingale.
\end{enumerate}
\end{studentexercise}
\begin{shortsolution}
\begin{enumerate}
  \item The payoff and paths stay fixed; only scenario weights change.
  \item One also needs \(\E^{\Pp}[Z_T]=1\); strict positivity gives
        equivalence.
  \item It is \(Z_t=\E^{\Pp}[Z_T\mid\cF_t]\).
  \item Zero drift gives a local martingale first; validity requires a
        true-martingale argument.
  \item The drift changes, while diffusion and quadratic variation remain.
  \item Prices discounted by the chosen numeraire are martingales.
\end{enumerate}
\end{shortsolution}

\subsection{The Black--Scholes sign drill}

\begin{studentexercise}{five lines from the historical model}
Assume constant \(r,\mu,\sigma\), with \(\sigma>0\), and
\(\dd S_t/S_t=\mu\dd t+\sigma\dd W_t^{\Pp}\). Write, in order, the
market price of risk, density process, \(\Q\)-Brownian motion,
\(\Q\)-stock dynamics, and discounted-stock dynamics.
\end{studentexercise}
\begin{shortsolution}
\[
\lambda=\frac{\mu-r}{\sigma},\qquad
Z_t=e^{-\lambda W_t^{\Pp}-\lambda^2t/2},\qquad
\frac{\dd\Q}{\dd\Pp}\bigg|_{\cF_T}=Z_T,
\]
\[
W_t^{\Q}=W_t^{\Pp}+\lambda t,\qquad
\frac{\dd S_t}{S_t}=r\dd t+\sigma\dd W_t^{\Q},\qquad
\boxed{\dd(e^{-rt}S_t)=\sigma e^{-rt}S_t\dd W_t^{\Q}.}
\]
The last expression is a stochastic integral and hence a local martingale;
calling it a true martingale needs an integrability argument.
\end{shortsolution}

\clearpage
\subsection{End-of-lecture checklist}

\begin{keyidea}{Nine checks for a change-of-measure calculation}
\begin{enumerate}
  \item State the original measure, filtration, and horizon.
  \item Identify the drift that must change and the diffusion that cannot.
  \item Write the candidate stochastic exponential with its compensator.
  \item Check nonnegativity, unit expectation, and equivalence.
  \item Use Novikov, boundedness, localization, or another valid criterion.
  \item State the new Brownian motion with the correct sign.
  \item Solve for the old differential before substituting into the SDE.
  \item When conditioning, use \(Z_t\), not the terminal density alone.
  \item For pricing, name the numeraire and discount every traded asset.
\end{enumerate}
\end{keyidea}

\begin{keyidea}{Final synthesis}
\begin{itemize}
  \item A density reweights unchanged scenarios, and Bayes' rule requires
        conditional renormalization.
  \item Gaussian tilting changes a mean without changing variance; Girsanov
        is its pathwise Brownian analogue.
  \item The half-quadratic-variation term makes a stochastic exponential a
        local martingale, not automatically a true martingale.
  \item Girsanov moves drift into the probability measure while preserving
        diffusion and quadratic variation.
  \item The negative market-price-of-risk density replaces \(\mu\) by \(r\)
        and makes the discounted stock a local martingale. Lecture 7
        develops the integrability, replication, and completeness steps
        needed for general valuation.
\end{itemize}
\end{keyidea}

\end{document}
